\documentclass[12pt]{article}
\usepackage[utf8]{inputenc}
\usepackage[T1]{fontenc}
\usepackage[margin=1in]{geometry}
\usepackage{setspace}
\usepackage{xcolor}
\usepackage[colorlinks=true, urlcolor=blue!60!black, linkcolor=blue!60!black]{hyperref}
\setlength{\parindent}{0em}
\setlength{\parskip}{0.6em}
\renewcommand{\baselinestretch}{1.0}
\begin{document}
\begin{center}
{\LARGE \textit{United States v. Nixon}, 418 U.S. 683 (1974)\par}
\vspace{0.5em}
{\large Case Brief --- Two-Page Edition\par}
\end{center}

\section{Memory Jogger}
The Watergate tapes case: presidential communications are presumptively privileged, but a generalized interest in confidentiality yields to the specific needs of a criminal trial (\href{https://supreme.justia.com/cases/federal/us/418/683/}{418 U.S. 683}).

\section{Detailed Case Facts}
After Nixon's secret Oval Office taping system was revealed and Archibald Cox was fired, Special Prosecutor Leon Jaworski, protected by a Justice Department regulation, obtained a Rule 17(c) subpoena for tapes of 64 conversations for the trial of seven Nixon aides, including Mitchell, Haldeman, and Ehrlichman. The grand jury named Nixon an unindicted coconspirator. Nixon moved to quash, claiming absolute privilege.

\section{Procedural History}
Chief Judge Sirica denied the motion and ordered \textit{in camera} production. The Supreme Court granted certiorari before judgment and ruled on July 24, 1974; Nixon resigned on August 9.

\section{Judicial Votes}
Unanimous, 8--0: Burger wrote for the Court; Rehnquist did not participate.

\section{Holding}
The order was affirmed. The appeal was proper, the dispute justiciable, and the subpoena valid under Rule 17(c). Presidential communications carry a presumptive constitutional privilege, but absent military or diplomatic secrets a generalized claim of confidentiality must yield to the specific need for evidence in a pending criminal trial.

\section{Analysis of Opinions}
Forcing a President into contempt to obtain review would be ``unseemly.'' Because the Executive was bound by its own regulation under \textit{Accardi}, the case presented ``concrete adverseness.'' Reaffirming \textit{Marbury}, Burger held that courts decide the scope of privilege; he grounded the privilege in Article II but, invoking Justice Jackson's \textit{Youngstown} concurrence, rejected absolute privilege. Balanced against the Fifth and Sixth Amendments and the public's right to evidence, occasional \textit{in camera} disclosure would not chill candid advice.

\section{Examples: Future Applications}
\textbf{Same-Side:} a special counsel's specific subpoena for recordings of Oval Office meetings about a bribe is enforceable. \textbf{Opposite-Side:} a subpoena for classified war plans meets ``utmost deference'' to military secrets. \textbf{Fence-Sitter:} a congressional subpoena for White House memoranda, which \textit{Nixon} left open.

\section{Critique}
Originalists such as Raoul Berger and Saikrishna Prakash doubt any constitutional privilege exists, and unitary-executive theorists dispute justiciability. William Van Alstyne observed that the Court defeated Nixon while giving later presidents a powerful privilege, and commentators note tension with \textit{Trump v. United States}.

\section{Key Quotations}
\begin{quote}
``A President and those who assist him must be free to explore alternatives in the process of shaping policies and making decisions, and to do so in a way many would be unwilling to express except privately. These are the considerations justifying a presumptive privilege for Presidential communications. The privilege is fundamental to the operation of Government, and inextricably rooted in the separation of powers under the Constitution.''
\end{quote}
\noindent --- Chief Justice Burger, Opinion of the Court, 418 U.S.\ at 708.

\end{document}
